\documentclass[10pt,a4paper]{amsart}
\usepackage[margin=19mm]{geometry}
\usepackage{amsmath,amssymb,mathtools}
\usepackage[scr=rsfs,cal=euler]{mathalfa}
\usepackage{xfrac}
\usepackage[unicode,hidelinks,bookmarks=false]{hyperref}
\newcommand{\ie}{i.e.\ }
\newcommand{\cf}{cf.\ }
\newcommand{\resp}{respectively\ }
\newcommand{\Subsection}{\subsection}
\providecommand{\nobreakdash}{\nobreak\hbox{-}}
\setlength{\emergencystretch}{2em}
\multlinegap0pt
\usepackage{bm}
\usepackage{bbm}
\usepackage{dsfont}
\usepackage{xspace}
\usepackage{cancel}
\usepackage{mathtools}
\usepackage{amsthm}
\makeatletter
\@ifundefined{thm}{\newtheorem{thm}{Theorem}}{}
\makeatother
\DeclareMathOperator{\Lin}{Lin}
\newcommand{\Aop}{\mathbb{A}}
\newcommand{\C}{\mathbb C}
\newcommand{\D}{\mathrm{D}}
\newcommand{\Proj}{\mathbb{P}}
\newcommand{\R}{\mathbb R}
\newcommand{\abs}[1]{\left\lvert#1\right\rvert}
\newcommand{\ball}{\mathrm{B}}
\newcommand{\dx}[1]{\,\mathrm{d}#1}
\newcommand{\hold}{\operatorname{C}}
\newcommand{\lebe}{\operatorname{L}}
\newcommand{\rfrak}{\mathfrak{r}}
\newcommand{\sobo}{\operatorname{W}}
\renewcommand{\L}{\mathrm{L}}
\title[On $\mathrm{BV}^{\mathbb{A}}$-Minimisers in two Dimensions]{On $\mathrm{BV}^{\mathbb{A}}$-Minimisers in two Dimensions}
\author{Ferdinand Eitler, Peter Lewintan}
\date{Source: arXiv:2508.10508v1 (CC BY 4.0)}
\begin{document}
\maketitle

\noindent\emph{Source and licence.} On $\mathrm{BV}^{\mathbb{A}}$-Minimisers in two Dimensions. Version arXiv:2508.10508v1 (CC BY 4.0), distributed under the Creative Commons Attribution 4.0 International licence, as recorded in the arXiv licence metadata for this exact version and shown on the abstract page https://arxiv.org/abs/2508.10508v1. Full rights notice: licenses/arxiv-2508.10508-CC-BY-4.0.txt.\par\smallskip

\noindent\textit{Editorial scope.} Counted unit: the Thm whose source label is \texttt{eq:KornType2\_Orlicz} in the article (source0.tex lines 969--975, source label \texttt{eq:KornType2\_Orlicz}), delivered together with the article's own complete original proof of it (source0.tex lines 976--990, one \texttt{proof} environment of 15 lines). No mathematics is written, repaired or re-derived: statement and proof are the article's own text line for line. Required context retained uncounted: eq:DifferentialOperator (lines 118--120); eq:CianchiAssumptionYoungFunctions (lines 909--918); eq:Korn\_Type\_Orlicz\_Paul\_Norm (lines 921--923); eq:Poincare\_Projection (lines 962--964). Every definition, display and label of those lines is retained unchanged. Omitted from this package and not redistributed: 1--117, 121--908, 919--920, 924--961, 965--968, 991--1261. Nothing inside a retained excerpt is omitted or altered: every retained body is delivered exactly as the article's lines are written, so no omission receipt of any kind accompanies this package. Citations: the retained excerpts carry no citation command at all, so no bibliography entry of the article is transcribed and no reference is redistributed. Reference adaptation: none was needed and none was made; every reference inside the delivered unit resolves inside it, so no unresolved reference or citation is produced. Printed slips kept literal: no printed word, symbol, hypothesis or number of the counted statement or of its proof was altered. Layout and class: the article's own class is \texttt{amsart} with packages including amsmath, amsthm, amssymb, enumerate, mathtools, stmaryrd, xcolor, bbm, libertine, euscript, mathrsfs, geometry, hyperref, cleveref; the wrapper is the lane's pinned \texttt{amsart} editorial wrapper at 10pt on A4, which restates the theorem-like environments the retained text needs and adds the article's own macro definitions that the retained text needs (\texttt{\textbackslash Aop}, \texttt{\textbackslash C}, \texttt{\textbackslash D}, \texttt{\textbackslash L}, \texttt{\textbackslash Lin}, \texttt{\textbackslash Proj}, \texttt{\textbackslash R}, \texttt{\textbackslash abs}, \texttt{\textbackslash ball}, \texttt{\textbackslash dx}, \texttt{\textbackslash hold}, \texttt{\textbackslash lebe}, \texttt{\textbackslash rfrak}, \texttt{\textbackslash sobo}), with extra wrapper packages (bm, bbm, dsfont, xspace, cancel, mathtools). The delivered numbering is the unit's own. Assets: no retained span contains a figure environment, a graphics-inclusion command, a PDF-page inclusion, a table or a TikZ picture; no image file is copied into this package, the package declares no asset dependency and no image file is redistributed with it. Licence: version arXiv:2508.10508v1 of this article is distributed under the Creative Commons Attribution 4.0 International licence, as recorded in the arXiv licence metadata for this exact version and shown on the abstract page https://arxiv.org/abs/2508.10508v1. A full rights notice is delivered at licenses/arxiv-2508.10508-CC-BY-4.0.txt. Characters: every retained excerpt is pure ASCII, so the delivered engine, XeLaTeX with its default Latin Modern text font, reports no missing character.
\begin{align}\label{eq:DifferentialOperator}
\Aop u\coloneqq \sum_{\abs{\alpha}=\ell}\Aop_{\alpha}\partial^\alpha u, \quad u\colon\R^n\to V,
\end{align}

\begin{thm}[Korn-type inequality in Orlicz spaces]
	Let $\Omega\subset\R^n$, $n\geq 2$ be an open and bounded set and $\Aop$ be a homogeneous, first-order elliptic differential operator of the form \eqref{eq:DifferentialOperator}. Moreover, let $\Psi$ and $\Phi$ be Young functions such that 
	\begin{equation}\label{eq:CianchiAssumptionYoungFunctions}
		t \int_0^t \frac{\Phi(s)}{s^2}\dx{s} \leq \Psi(ct) \quad\mbox{and}\quad t\int_0^t \frac{\Psi^\ast(s)}{s^2}\dx{s}\leq \Phi^\ast(ct),
	\end{equation}
	hold for all $t\geq 0$ and some $c>0$. Then there exist a constant $C>0$ such that 
	\begin{equation}\label{eq:Korn_Type_Orlicz_Paul}
		\int_\Omega \Phi(\abs{\D u})\dx{x} \lesssim \int_\Omega \Psi(C\abs{\Aop u})\dx{x}\quad\mbox{for all}\quad u\in \sobo^{1,\Psi}_0(\Omega;V).
	\end{equation}
\end{thm}

\begin{equation}\label{eq:Korn_Type_Orlicz_Paul_Norm}
	\|\D u \|_{\lebe^\Phi(\Omega;\Lin(\R^n;V))} \leq C\, \|\Aop u\|_{\lebe^\Psi(\Omega;W)}.
\end{equation}

	\begin{align}\label{eq:Poincare_Projection}
		\|u-\Pi^\ball_{\Aop} u\|_{\L^\Psi(\ball;V)} \leq c\, \omega_n\, r \|\Aop u\|_{\L^\Psi(\ball;W)}, 
	\end{align}

\begin{thm}[Korn-type inequality]
	Let $r>0$ as well as $\Psi$ and $\Phi$ be Young functions such that \eqref{eq:CianchiAssumptionYoungFunctions} is satisfied. Then for each $\beta>0$ there exists a constant $c=c(\Phi,\Psi,n,r,\Aop)>0$ such that 
	\begin{align}\label{eq:KornType2_Orlicz}
		\inf_{\rfrak\in\ker(\Aop)} \|\D (u-\rfrak)\|_{\L^\Phi(\ball_r(x_0);\Lin(\R^n;V))}\leq c \bigg(1+\frac{1}{r}\bigg)\|\Aop u \|_{\L^\Psi(\ball_{2r}(x_0);W)}
	\end{align}
	for all $u\in\sobo^{\Aop,\Psi}(\ball_{2r}(x_0))$, where $\Aop$ is a  $\C$-elliptic differential operator of the form \eqref{eq:DifferentialOperator}.
\end{thm}

\begin{proof}
	Let $\varrho\in\hold^\infty_c(\R^n;[0,1])$ be a cut-off function satisfying $\chi_{\ball_{r}(x_0)}\leq \varrho\leq \chi_{\ball_{2r}(x_0)}$ and $\abs{\nabla\varrho}\leq\frac{2}{r}$. Moreover, we recall from \eqref{eq:Poincare_Projection} that the  Poincar\'e-type inequality 
	\begin{equation}\label{eq:Poincare_KornProof}
		\|u-\Pi_\Aop^\ball u\|_{\L^\Psi(\ball_{2r}(x_0);V)} \leq c\,\omega_n\, r\|\Aop u\|_{\L^\Psi(\ball_{2r}(x_0);W)},
	\end{equation}
	holds true. Therefore, taking into account \eqref{eq:Korn_Type_Orlicz_Paul_Norm} on $\Omega=\ball_{5r}(x_0)$ applied for the function $\varrho(u-\Proj_\ball u)\in\sobo^{1,\Psi}_0(\ball_{5r}(x_0);V)$ in conjunction with \eqref{eq:Poincare_KornProof}, leads to 
	\begin{align*}
		\|\D (u-\Pi_\Aop^\ball u)\|_{\L^\Phi(\ball_{r}(x_0))} 
		&\leq \|\D (\varrho(u-\Pi_\Aop^\ball u))\|_{\L^\Phi(\ball_{5r}(x_0)}\\
		& \leq c \,\|\Aop (\varrho(u-\Pi_\Aop^\ball u))\|_{\L^\Psi(\ball_{5r}(x_0);W)}\\
		&\leq c \,\bigg( \frac1r\left\|u-\Pi_\Aop^\ball u\right\|_{\L^\Psi(\ball_{2r}(x_0);V)} +  \|\Aop u\|_{\L^\Psi(\ball_{2r}(x_0); W)}\bigg)\\
		&\leq c\,\bigg(1+\frac{1}{r}\bigg) \|\Aop u\|_{\L^\Psi(\ball_{2r}(x_0);W)}.
	\end{align*}
	This yields the desired inequality \eqref{eq:KornType2_Orlicz}, as we can always bound the left-hand side by the infimum taken over all $\rfrak\in\ker(\Aop)$, which finishes the proof. 
\end{proof}

\end{document}
